\documentclass[border=1cm]{standalone}
\usepackage{chronos}
\begin{document}
		\begin{chronos}
			[
			date centric,
			timeline={
				start date=2005,
				end date=2012,
			},
			]% La principal dificultad al resolver ecuaciones diferenciales singulares es el
    % comportamiento singular que ocurre en \(r = 0\). Sin embargo, se han utilizado
    % muchos métodos para resolver este tipo de ecuaciones diferenciales singulares.
    % Wazwaz (2005)
    % Khan (2012) aplicaron el método de transformación diferencial (DTM)
    % Chowdhury y Hashim (2009) utilizaron el método de perturbación por homotopía (HPM).
    % Baranwal \textit{et al}. (2012) aplicaron un enfoque híbrido
    % Parand \textit{et al}. (2009) utilizaron una técnica seudoespectral basada en las funciones de Legendre
    % racionales y la integración de Gauss--Radau para la solución de la ecuación de
    % Lane--Emden de tipo (1).
			\chronosevent {date={2005},
				yshift=5pt,
				special date={2000},
				name=C,
				name content=Wazwaz }
			\chronosevent {date={2009},		
				name=B,
				name content=Chowdhury y Hashim,
				yshift=5pt}
            \chronosevent {date={2009},
				as is,
				name=B,
				name content=Parand \textit{et al},
				yshift=30pt}
			\chronosevent {date={2012},
				as is,
				name=A,
				name content=Khan,
				yshift=5pt}
            \chronosevent {date={2012},
				as is,
				name=A,
				name content=Baranwal \textit{et al},
				yshift=5pt}
			\chronosmaintitle {name=Theories of X, at=current bounding box.north west,yshift=30pt,anchor=north west}
		\end{chronos}
\end{document}